\documentclass[parskip=full]{scrartcl}

\usepackage[ngerman,main=english]{babel}
\usepackage[colorlinks,allcolors=blue]{hyperref}
\usepackage{refenums}
% refenums references with zref-clever; capitalize its names of other references, too (e.g., "Section 3")
\zcsetup{cap}

%Always add a dot behind the definition
\renewcommand{\refenumenclosing}[1]{#1.}

% \bexample ... \eexample runs the code in between and then prints its source lines.
% \bsetup{<explanation>} ... \esetup does the same for setup code, which has no output; the explanation goes into the box.
% Plain commands, not environments: a group would discard the macros \setupRefEnums defines.
\usepackage[listings]{tcolorbox}
\usepackage{booktabs}
\usepackage{array}
\tcbset{
  demo listing/.style={
    listing file=\jobname.tex,
    left=8mm,
    listing options={
      language={[LaTeX]TeX},
      basicstyle=\small\ttfamily,
      columns=fullflexible,
      breaklines,
      breakatwhitespace,
      firstline=\the\bexampleline,
      lastline=\the\eexampleline,
      firstnumber=\the\bexampleline,
      numbers=left,
      numberstyle=\tiny,
      numbersep=8pt
    }
  }
}
\newcount\bexampleline
\newcount\eexampleline
\newcommand{\bexample}{\global\bexampleline=\numexpr\inputlineno+1\relax}
\newcommand{\eexample}{%
  \global\eexampleline=\numexpr\inputlineno-1\relax
  \tcbinputlisting{demo listing, listing only, colback=green!5!white, colframe=green!25}%
}
\newcommand{\bsetup}[1]{\gdef\setupexplanation{#1}\bexample}
\newcommand{\esetup}{%
  \global\eexampleline=\numexpr\inputlineno-1\relax
  \tcbinputlisting{demo listing, comment and listing, comment={\setupexplanation}, colback=orange!5!white, colframe=orange!40}%
}

\begin{document}

\title{Demonstration of the \texttt{refenums} package}
\date{}
\author{}
\maketitle

\vspace{-12ex}
This document demonstrates the \href{https://ctan.org/pkg/refenums}{\texttt{refenums}} package with the development documents of a reference manager.
The content is a simplified version of \href{https://github.com/JabRef/jabref/tree/main/docs}{JabRef's} requirements and decision records, written in \href{https://adr.github.io/madr/}{MADR}.

Each section starts with an orange box setting up a kind of item; this produces no output.
Each green box shows the \LaTeX{} code producing the output right above it.
Items are defined and referenced with shorthands such as \verb|\defReq| and \verb|\refReq|, which \verb|\setupRefEnum{Req}| creates for \verb|\defRefEnum{Req}| and \verb|\refEnum{Req}|.

\tableofcontents
\clearpage

\section{Features}
\label{sec:features}

\bsetup{Features print only their short form (e.g., ``Feature~1'') and each becomes a subsection.}
\setupRefEnum[separator={ }, enclosing=subsection]{Feature}{ONLYSHORT}
% With the default separator "-", this would print "Feature-1"
\esetup

\bexample
\defFeature{Save a library}{save}
As a researcher, I want to save my library to a \texttt{.bib} file, so that I can cite its entries in \LaTeX{}.
\eexample

\bexample
\defFeature{Autosave}{autosave}
As a researcher, I want every change to be saved automatically and a backup to be kept, so that I do not lose work after a crash.
\eexample

\section{Risks}
\label{sec:risks}

\bsetup{Risks use no separator (e.g., ``Risk~1 (R1)'').}
\setupRefEnum[separator={}]{R}{Risk}
% With the default separator "-", this would print "R-1"
\esetup

\bexample
\defR{Concurrent modification}{concurrent}
Another process, e.g., a second JabRef instance on a network share, writes the library file while \refFeature{save} is running.
\eexample

\bexample
\defR{Incomplete backup}{backup}
A failure while writing a backup during \refFeatureFull{autosave} replaces a valid backup.
\eexample

\section{Requirements}
\label{sec:requirements}

\bsetup{Requirements use the default separator (e.g., ``Requirement~1 (Req-1)'').}
\setupRefEnum{Req}{Requirement}
% Titles are set in italics; titleformat={} would set them upright
\esetup

\bexample
\defReq{Detect concurrent modification}{detect}
If the library file changed since the save started, the save is aborted with an error.
Mitigates \refRFull{concurrent}.
\eexample

\bexample
\defReq{Keep valid backups}{keep}
A failure while writing a backup leaves the previous backup untouched.
Mitigates \refRFull{backup}.
\eexample

\section{Decisions}
\label{sec:decisions}

\bsetup{Each decision record becomes a subsection. The options a record considers are defined inline, within the text. Both use the default separator.}
\setupRefEnum[enclosing=subsection]{ADR}{Decision Record}
\setupRefEnum[inline]{Opt}{Option}
% Without inline, \defOpt would print "Option 1 (Opt-1): ..."; \defOptInline prints the inline form for a single option
\esetup

\bexample
\defADR{Replace files atomically}{atomic}

\paragraph{Context and Problem Statement}
How do we write a file so that \refReqFull{detect} and \refReqFull{keep} hold?

\paragraph{Considered Options}
\begin{itemize}
\item \defOpt{Overwrite in place}{inplace}
\item \defOpt{Write a temporary file and move it over the target}{tmp}
\end{itemize}

\paragraph{Decision Outcome}
Chosen option: \refOptFullColon{tmp}, because a failed write leaves the original file untouched and the target can be checked for changes right before the move.
\eexample

\section{Milestones}
\label{sec:milestones}

\bsetup{Milestones combine print name and short name (e.g., ``Milestone~M1'').}
\setupRefEnum[combined]{M}{Milestone}
% Without combined, this would print "Milestone 1 (M-1)"
\esetup

\bexample
\defM{Safe saving}{safe}
Reached when \refReq{detect} and \refReq{keep} are implemented as decided in \refADR{atomic} and covered by unit tests.
This mitigates \zcref{enum:R:concurrent,enum:R:backup}.
\eexample

\section{German requirements}
\label{sec:german}

\bsetup{The plural of the print name defaults to the print name with an ``s'' appended (e.g., ``Requirements''). German needs its own plural. \texttt{\textbackslash zcref} follows the language of the text, here set with babel.}
\setupRefEnum[plural=Anforderungen]{Anf}{Anforderung}
% Without plural, \zcref would print "Anforderungs"
\esetup

\bexample
\selectlanguage{ngerman}
\defAnf{Bibliothek speichern}{speichern}
Die Bibliothek wird als \texttt{.bib}-Datei gespeichert.
\eexample

\bexample
\defAnf{Sicherung behalten}{sicherung}
Ein fehlgeschlagenes Schreiben lässt die vorherige Sicherung unverändert.
\eexample

\bexample
Diese Version erfüllt \zcref{enum:Anf:speichern,enum:Anf:sicherung}.
\selectlanguage{english}
\eexample

\section{References}
\label{sec:references}

The package \texttt{refenums} offers several ways to reference an item.
They differ in how much of the item they show: its kind, its identifier, its title, or a combination.
The following table shows each of them for \refReq{keep} from \zcref{sec:requirements}.

\begin{center}
\small
\renewcommand{\arraystretch}{1.3}
\setlength{\tabcolsep}{4pt}
\begin{tabular}{@{}l>{\footnotesize}ll@{}}
\toprule
Aim & Command & Output \\
\midrule
Identifier & \verb|\refReq{keep}| & \refReq{keep} \\
Identifier (title) & \verb|\refReqFull{keep}| & \refReqFull{keep} \\
Identifier: title & \verb|\refReqFullColon{keep}| & \refReqFullColon{keep} \\
Title & \verb|\refReqTitle{keep}| & \refReqTitle{keep} \\
\midrule
\multicolumn{3}{@{}l}{\emph{With the kind of the item}} \\
Identifier & \verb|\refReq*{keep}| & \refReq*{keep} \\
Identifier (title) & \verb|\refReqFull*{keep}| & \refReqFull*{keep} \\
Identifier: title & \verb|\refReqFullColon*{keep}| & \refReqFullColon*{keep} \\
\bottomrule
\end{tabular}
\end{center}

Each shorthand has a long form, which takes the kind of the item as first argument.
For instance, \verb|\refReqFull{keep}| is short for \verb|\refEnumFull{Req}{keep}|.
The command \verb|\refReq*{keep}| is also short for \verb|\zcref{enum:Req:keep}| of \href{https://ctan.org/pkg/zref-clever}{zref-clever}, which also references several items at once, as in \zcref{sec:milestones}.

The following list references every item of this document in the form ``Identifier (title)''.

\bexample
\begin{itemize}
\item Features: \refFeatureFull{save}, \refFeatureFull{autosave}
\item Risks: \refRFull{concurrent}, \refRFull{backup}
\item Requirements: \refReqFull{detect}, \refReqFull{keep}
\item Decision records: \refADRFull{atomic}
\item Options: \refOptFull{inplace}, \refOptFull{tmp}
\item Milestones: \refMFull{safe}
\end{itemize}
\eexample

\end{document}
